\section{Spectral data of a supersymmetric realization: \((\mathcal H_{\rm phys},H,\mathcal Q)\), supercharge and compactification in dimension \(11\)}
\label{sec:md-sucesora-terna-supercarga-m}

The interface dimensional toward theory M requires distinguish the exact structure
of split, the lift operatorial and the fields fisicos of the codomain. Let
\[
 \mathfrak S_M=(\mathcal H_{\mathrm{phys}},H,Q)
\]
a spectral triple in which \(\mathcal H_{\mathrm{phys}}\) is the Hilbert
space after imposing the gauge constraints, \(H\) is a nonnegative
self-adjoint operator, and \(Q\) is an odd, densely defined supercharge
compatible with the grading
\(\mathcal H_{\mathrm{phys}}=\mathcal H_+\oplus\mathcal H_-\). The relation
\begin{equation}
 Q:\mathcal H_+\rightleftarrows\mathcal H_-,
 \qquad
 Q^2=H
 \label{eq:md-sucesora-supercarga-cuadrado}
\end{equation}
is interpreted on its common domain. It does not replace the relativistic
supersymmetric algebra; it is its Hermitian reduction of dimension \(1\) and fixes the minimal datum
that a spectral realization must preserve.

\subsection{Operator lifting of the APP--TRIT--TPK state}

Let \(\mathcal H_{\mathrm{HMT}}\) be the completion of the surviving sections
with phase, orientation, quotient, carry, boundary, route, and
memory, and let
\[
 W:\mathcal H_{\mathrm{HMT}}\longrightarrow\mathcal H_{\mathrm{phys}}
\]
a realization isometry. An admissible operator lift preserves
the APP composition, the TRIT orientation, and the full TPK dynamics,
and intertwines
\begin{equation}
 W\widehat H=HW,
 \qquad
 W\widehat Q=QW,
 \qquad
 \widehat Q^{2}=\widehat H.
 \label{eq:md-sucesora-lift-app-supercarga}
\end{equation}
The first equality transports the generator, the second preserves the grading,
and the third prevents introducing \(H\) as a datum foreign to the genealogical
antecedent. These identities are obligations of the bridge: the existence of
fermionic and bosonic completions alone does not prove
\eqref{eq:md-sucesora-supercarga-cuadrado}.

\Needspace{42\baselineskip}
\subsection{Physical undecadimensional codomain}

The realization in dimension \(11\) is expressed by
\[
 \mathfrak R_M:\mathcal X_{\mathrm{HMT}}^{\mathrm{enr}}
 \longrightarrow
 \mathcal X_{11}^{\mathrm{phys}},
\]
whose codomain carries metric \(g_{MN}\), the \(3\)-form \(C_{MNP}\), the
gravitino \(\psi_M\), the action, and its supersymmetry transformations.
Compactification along one direction must simultaneously transport these
fields, the inner product, and the grading of \(\mathfrak S_M\). Strings
appear on the dimension-(10) screen; branes couple to the
\(3\)-form and its duals; and T-duality exchanges momentum and
winding through
\[
 (m,w,R)\longmapsto(w,m,R^{-1})
\]
without altering the spectrum of the compact sector.

The Polyakov action
\[
 S_{\mathrm P}[X,h]
 =-\frac{T}{2}\int d^2\sigma\,
 \sqrt{-h}\,h^{ab}g_{MN}(X)
 \partial_aX^M\partial_bX^N
\]
fixes the worldsheet dynamics once \(T\), \(h_{ab}\), the maps \(X^M\), and the boundary conditions have been determined. The screen chains \(12\to11\to10\) and \(26\to24\), the radius involution, and the lattice reduction are prior algebraic results; the action, spectrum, and amplitudes constitute the subsequent physical realization.

\subsection{Compatibility of the spectral triple \((\mathcal H_{\rm phys},H,\mathcal Q)\) with the physical codomain in dimension \(11\)}

A realization is admissible if it preserves at the same time:
\begin{enumerate}
 \item the compatible prolongation of the state to every depth;
 \item the domains and codomains of \(H\), \(Q\), and the observables;
 \item the intertwining identities
       \eqref{eq:md-sucesora-lift-app-supercarga};
 \item the unitary action of T-duality on the spectrum;
 \item the descent of the fields of dimension \(11\) to the stated
       compactification;
 \item compatibility with the physical inner product and the gauge
       constraints.
\end{enumerate}
The failure of any of these conditions distinguishes an analogy of ranks
from a physical realization. The formulation does not automatically identify the
exceptional corridor with M-theory: both publications follow from the same
enriched state, but use distinct functors and codomains.
